\documentclass[10pt,a4paper]{article}
\usepackage[T1]{fontenc}
\usepackage[utf8]{inputenc}
\usepackage[spanish,es-nodecimaldot]{babel}
\usepackage[margin=1.7cm,top=1.2cm,bottom=1.0cm]{geometry}
\usepackage{xcolor}
\usepackage{titlesec}
\usepackage{enumitem}
\usepackage{parskip}
\setlength{\parskip}{3.2pt plus 1pt minus 0.5pt}
\usepackage{tcolorbox}
\usepackage{multicol}
\usepackage[hidelinks]{hyperref}

\definecolor{azul}{HTML}{123B5E}
\definecolor{gris}{HTML}{4A4A4A}
\definecolor{fondo}{HTML}{F2F5F8}

\titleformat{\section}{\normalfont\large\bfseries\color{azul}}{}{0pt}{}
\titlespacing*{\section}{0pt}{10pt}{4pt}

\setlist[itemize]{leftmargin=1.1em,itemsep=1.5pt,topsep=2pt,parsep=0pt}

\pagestyle{empty}

\input{ficha-version}

\newcommand{\barra}{\textcolor{azul}{\rule{\linewidth}{1.2pt}}}

\begin{document}

{\color{azul}\Huge\bfseries MuniANCI}\\[2pt]
{\color{gris}\large Cumplimiento de la Ley 21.663 para organismos del Estado, sin conexión a
internet}

\vspace{4pt}
\barra

\vspace{6pt}

La Ley 21.663 obliga a los organismos de la Administración del Estado a aplicar medidas
permanentes de ciberseguridad (Art. 7\textdegree) y a reportar los incidentes de efecto
significativo al CSIRT Nacional en plazos perentorios: alerta temprana en 3 horas,
actualización en 72, informe final en 15 días corridos (Art. 9\textdegree). MuniANCI mide
el estado real de la institución frente a esas obligaciones y produce la evidencia, en un
solo ejecutable que no necesita servidor, base de datos ni salida a internet.

\begin{multicols}{2}

\section{Qué hace}

\textbf{Escanea la red institucional.} Descubrimiento nativo por ARP e ICMP, inventario de
equipos y servicios, y correlación de vulnerabilidades contra una base CVE que viaja dentro
del producto. Prioriza con el catálogo KEV de CISA, que distingue lo que se está explotando
hoy de lo que solo figura en una lista.

\textbf{Evalúa el cumplimiento artículo por artículo.} Con una distinción que casi nadie
hace: lo exigible a la institución se mide como cumplimiento, mientras que el
Art. 8\textdegree, que obliga a los operadores de importancia vital, se mide como madurez
voluntaria y se rotula como tal.

\textbf{Entrega el expediente completo.} Informe técnico por dominio, informe ejecutivo de
una página, plan de remediación priorizado en formato OSCAL, y el JSON para el CSIRT con la
taxonomía de la Res. Ex. N\textdegree 7/2025 transcrita del Diario Oficial.

\textbf{Mide la evolución.} Histórico por institución y deriva por control entre escaneos:
una brecha nueva, persistente, resuelta o reaparecida. La reaparecida es la que habla del
proceso interno.

\textbf{Responde preguntas normativas.} Asistente legal que trabaja sobre el texto de las
normas y cita el documento del que sacó cada respuesta. Si no puede respaldar una
referencia con el texto a la vista, no la entrega.

\section{Por qué offline}

Nada institucional sale del equipo. No hay nube, no hay telemetría, no hay cuenta que
crear. El modelo de lenguaje, el índice de vulnerabilidades y el corpus normativo viajan
dentro del instalador y funcionan en un equipo aislado de la red.

Es una decisión de diseño: un producto de cumplimiento que exporta el inventario de
vulnerabilidades de una institución a un servidor de terceros traslada el riesgo en vez de
reducirlo.

\section{Multi-marco}

Además de la Ley 21.663, mide el \textbf{Decreto 7 de 2023} de MINSEGPRES, cuyas cinco
funciones del Título Tercero coinciden con las del NIST CSF, y sitúa a la institución en la
fase de la \textbf{Ley 21.180} que le corresponde.

\end{multicols}

\vspace{-4pt}

\begin{tcolorbox}[colback=fondo,colframe=azul,boxrule=0.6pt,arc=2pt,left=6pt,right=6pt,
top=5pt,bottom=5pt]
{\bfseries\color{azul}Lo que este producto no hace}

\begin{itemize}
\item \textbf{No firma electrónicamente.} El paquete de evidencia lleva un manifiesto
SHA256 que se verifica con \texttt{certutil} o \texttt{Get-FileHash}, herramientas que
Windows ya trae. Es verificación de integridad. Bajo la Ley 19.799, solo una firma
electrónica avanzada de un prestador acreditado da a un documento de un órgano del Estado
la calidad de instrumento público.
\item \textbf{No entrega asesoría legal.} Cita los artículos y transcribe las categorías
oficiales; la calificación jurídica de un incumplimiento concreto la valida un abogado.
\item \textbf{No clasifica a la institución por su cuenta.} La condición de operador de
importancia vital nace de una resolución fundada de la Agencia (Arts. 5\textdegree\ y
6\textdegree); el producto no la supone.
\end{itemize}
\end{tcolorbox}

\vspace{2pt}

\begin{multicols}{2}
\section{Requisitos}
\begin{itemize}
\item Windows 10 o superior, 64 bits.
\item 16 GB de RAM recomendados para el Asistente; el escáner funciona con menos.
\item Sin privilegios de administrador para el escaneo de red ni para programar el
reescaneo periódico.
\item Sin conexión a internet en ningún momento.
\end{itemize}

\section{Configuración local}

La identidad de la institución, los plazos, el alcance de red, la retención del histórico y
el contenido del informe se ajustan en un archivo de texto junto al ejecutable, que se abre
con el Bloc de notas. El panel de ajustes de la aplicación hace lo mismo sin tocarlo.

\section{Desarrollo sin costo para las FF.AA. y las policías}

El desarrollo no se le cobra a las Fuerzas Armadas, a Carabineros de Chile ni a la Policía
de Investigaciones: correcciones, adaptaciones al procedimiento interno y requerimientos
nuevos entran sin costo. El licenciamiento se acuerda por separado.
\end{multicols}

\vspace{4pt}
\barra

\vspace{2pt}

{\footnotesize\color{gris}
\textbf{Felipe Carvajal Brown} \quad\textbullet\quad
\href{mailto:fcarvajalbrown@gmail.com}{fcarvajalbrown@gmail.com} \quad\textbullet\quad
MuniANCI versión \version

\vspace{2pt}
Las referencias legales de esta ficha provienen del texto oficial de la Ley 21.663
publicado por la Biblioteca del Congreso Nacional y de las resoluciones de la ANCI
publicadas en el Diario Oficial. Este documento es material informativo y no constituye
asesoría legal.
}

\end{document}
